% Part: methods
% Chapter: proofs
% Section: starting-proofs

\documentclass[../../../include/open-logic-section]{subfiles}

\begin{document}

\olfileid[pt]{mth}{prf}{str}

\olsection{Começando uma Prova}

Mas por onde começar?

Foi-lhe dado algo a provar, de modo que isso deveria ser a última
coisa mencionada na prova (você pode, obviamente, \emph{anunciar} que
vai prová-lo no começo, mas não convém usá-lo como suposição). Escreva
o que você está tentando provar no rodapé de uma folha de papel em
branco---assim você não perde de vista seu objetivo.

Em seguida, você pode ter algumas suposições que pode usar (isso
ficará mais claro quando falarmos sobre o \emph{tipo} de prova que você
está fazendo na próxima seção). Escreva-as no topo da página e
certifique-se de sinalizar que são suposições (isto é, se você está
supondo $p$, escreva ``suponha que $p$'' ou ``assuma que $p$''). Por
fim, pode haver na questão algumas definições que você precisa
conhecer. Pode ser que lhe digam para usar uma definição específica, ou
pode haver várias definições nas suposições ou na conclusão para a qual
você está trabalhando. \emph{Anote-as e certifique-se de entender o que
significam.}

A forma como você organiza sua prova também dependerá da forma da
questão. A próxima seção fornece detalhes sobre como organizar sua
prova com base no tipo de sentença.

\end{document}
